% Generated draft: requires mathematical and source review.
\documentclass[11pt]{article}
\usepackage[margin=25mm]{geometry}
\usepackage{fontspec}
\setmainfont{Latin Modern Roman}
\usepackage{amsmath,amssymb,mathtools,unicode-math}
\setmathfont{Latin Modern Math}
\usepackage{enumitem,needspace,xurl,hyperref}
\setlength{\parindent}{0pt}
\setlength{\parskip}{5pt}
\setlength{\emergencystretch}{3em}
\newcommand{\N}{\mathbb{N}}
\newcommand{\Z}{\mathbb{Z}}
\newcommand{\Q}{\mathbb{Q}}
\newcommand{\R}{\mathbb{R}}
\newcommand{\C}{\mathbb{C}}
\newcommand{\F}{\mathbb{F}}
\newcommand{\A}{\mathbb{A}}
\newcommand{\PP}{\mathbb{P}}
\newcommand{\E}{\mathbb{E}}
\newcommand{\eps}{\varepsilon}
\newcommand{\dd}{\,\mathrm{d}}
\begin{document}
\section*{UnsolvedMath ID: 20003355}
\textbf{Source problem code:} catalogue entry
\subsection{Definitions and mathematical statement}
Question 53 (L. Silberman) asks: give an ergodic-theoretic proof of equidistribution of special points on $Y=\operatorname{SL}_2(\Z)\backslash\mathbb H$, the modular surface. The aid identifies these as CM packets. A precise fundamental-discriminant version is the following: for each negative fundamental discriminant $D<-4$, let $\mathcal Q_D$ be the finite set of $\operatorname{SL}_2(\Z)$-classes of primitive positive-definite integral binary quadratic forms $[a,b,c]$ with $b^2-4ac=D$, and put
\[
 z_{[a,b,c]}=\frac{-b+\mathrm i\sqrt{|D|}}{2a},\quad
 \mu_D=\frac1{|\mathcal Q_D|}\sum_{[a,b,c]\in\mathcal Q_D}\delta_{[z_{[a,b,c]}]},
 \qquad m_Y=\frac3\pi\frac{dx\,dy}{y^2}.
\]
The requested conclusion is $\mu_D\to m_Y$ weakly as $D\to-\infty$ through all negative fundamental discriminants, proved by ergodic methods. Weak probability convergence includes absence of escape to the cusp. The one-sentence source does not itself specify fundamental versus nonfundamental order discriminants or packet weights, so this is an explicit version selected from the aid's discussion, not a complete reconstruction of all intended variants. The aid reports the full Duke equidistribution theorem as analytically proved, while locating purely ergodic proofs only under fixed splitting conditions; this is a supplied literature qualification, not an independently checked impossibility claim.

Its companion dynamical closure criterion is conditional. Let $X$ be locally compact and metrizable, $T:X\to X$ a homeomorphism, $\eta_j$ probability measures, $m_X$ a target probability, and $\pi:X\to Y$ continuous with $\pi_*m_X=m_Y$. Suppose the $\eta_j$ are tight and asymptotically $T$-invariant, so every weak probability limit is $T$-invariant; suppose every such limit has maximal measure-theoretic entropy, and maximal-entropy rigidity identifies every such invariant probability with $m_X$. Then $\eta_j\to m_X$, and consequently $\pi_*\eta_j\to m_Y$. The required lift, tightness, invariance, entropy and rigidity inputs must be established for the actual packets; none is supplied as an unconditional construction for all discriminants here.

The finite-place obstruction reported in the aid is
\[
 \forall P\subset\{\text{primes}\}\text{ finite}\ \forall B>0\ 
 \exists\text{ negative fundamental }D<-B:\quad
 \left(\frac Dp\right)=-1\quad\text{for every }p\in P,
\]
where the symbol is the Kronecker symbol, including at $p=2$. Thus all primes in any fixed finite menu are simultaneously inert in infinitely many imaginary quadratic fields. A method requiring a split prime from that fixed menu misses infinitely many discriminants. This does not rule out varying-place methods or other ergodic proofs; the aid describes a CRT-plus-Dirichlet argument for this obstruction.

The remaining extracted text is a separate Section 10, ``Arithmeticity'', preceded by a page header. It sets $G=\operatorname{SL}_2(\R)^k$ with $k\geq2$, takes $U^+$ and $U^-$ to be the products of upper and lower unipotent subgroups, and takes lattices $\Gamma^\pm\subset U^\pm$. Its specified ``irreducibility'' means that every coordinate projection $G\to\operatorname{SL}_2(\R)$ is injective on each $\Gamma^\pm$. It introduces a following conjecture communicated by H. Oh but gives no conjecture text. These assumptions are preserved as neighboring context, not attached to Silberman's CM-equidistribution question.
\subsection{Short English statement}
Silberman asks for an ergodic proof of CM-packet equidistribution on the modular surface. The aid provides sufficient tightness and entropy inputs and explains why any fixed finite menu of split-prime actions cannot cover every fundamental discriminant. The appended arithmeticity setup belongs to a different, missing conjecture.
\subsection{Sources}
\begin{itemize}
\item Local UnsolvedMath registry, permanent ID 20003355, exact source code catalogue entry; \texttt{aim-other-notes.json}, notes[161], Question 53, L. Silberman.
\item American Institute of Mathematics, \emph{Emerging applications of measure rigidity}, \url{https://aimath.org/WWN/measrigid/measrigid.pdf}. The supplied verification identifies an OCR boundary between Question 53, a page header, and Section 10 on Arithmeticity; its claim of PDF inspection is source metadata, not an inspection performed here. H. Oh is credited only for the following arithmeticity conjecture.
\item Supplied research attempt 1: reduction, problem status partially\_solved, rigor label proved, novelty confidence low. It reports Duke's analytic theorem, fixed-splitting ergodic results, the tightness--entropy--rigidity criterion and finite-place obstruction. Machine-generated, not peer reviewed; independently verify before citation. No complete publication references for Duke or the ergodic results are supplied; original remarks and literature fields are empty.
\end{itemize}
\end{document}
